\originalchapter{课程作业: 势论与波能量}
\originalsection{作业4: 2011年10月6日交}
原题 I--V 是商业教材题号引用, 见缺项表.
\begin{exercise}\label{ps:4VI}
\textbf{原题 VI.} 单位球中 $\Delta u=f$, $u|_{\partial B}=g$, 证明存在与 $f,g$ 无关的 $C$ 使 $\|u\|_\infty\le C(\|f\|_\infty+\|g\|_\infty)$. 补充 $f$ 有界、$u$ 有连续边界迹, 以 $\sup$ 解释范数.
\end{exercise}
\begin{solution}
置 $M=\|g\|_\infty$, $F=\|f\|_\infty$, $b=(1-|x|^2)/(2n)$, 则 $b\ge0$, $\Delta b=-1$. 对 $v=u-M-Fb$ 有 $\Delta v=f+F\ge0$, 边界 $v\le0$, 最大原理给 $u\le M+Fb$. 同样用于 $-u$ 得 $|u|\le M+F/(2n)$. 可取 $C=\max(1,1/(2n))=1$. 开球光滑并不自动使 $f$ 有界, 此条件不可省.
\end{solution}
\originalsection{作业5: 2011年10月13日交}
本节沿原作业的符号 $\Delta_yG(x,y)=\delta_x$, 故 $G$ 为负; 与正文的 $-\Delta$ 正 Green 核相反.
\begin{exercise}\label{ps:5I}
\textbf{原题 I.} 三维单位球的核
\[
G(x,y)=-\frac1{4\pi|x-y|}+\frac1{4\pi|x||x/|x|^2-y|},\quad x\ne0,
\]
以及 $G(0,y)=-1/(4\pi|y|)+1/(4\pi)$, 直接证明 $G\le0$.
\end{exercise}
\begin{solution}
计算 $|x|^2|x/|x|^2-y|^2-|x-y|^2=(1-|x|^2)(1-|y|^2)>0$. 所以像点项的分母较大, 正项小于负项的绝对值, $G<0$ 于 $x\ne y$. $x=0$ 时 $|y|<1$ 直接得到负性. 对角线上核发散为 $-\infty$, 不作有限点值断言.
\end{solution}
\begin{exercise}\label{ps:5II}
\textbf{原题 II.} 计算 $\int_B G(x,y)\dd y$, 推出 $\int_B|G|\le1/6$, 并从表示式证明 $\Delta u=f$, 边值 $g$ 的一致数据估计.
\end{exercise}
\begin{solution}
$w(x)=(|x|^2-1)/6$ 满足 $\Delta w=1$, 边值0. Green 表示式给 $w(x)=\int_BG(x,y)\dd y$. 由负性,
$\int_B|G|=(1-|x|^2)/6\le1/6$. 又 $P(x,\sigma)=\partial_{\nu_\sigma}G=(1-|x|^2)/(4\pi|x-\sigma|^3)\ge0$. 将常函数1代入表示式得 $\int_{\partial B}P=1$. 因此
$|u(x)|\le\|f\|_\infty/6+\|g\|_\infty$. 仍须 $f$ 有界且边界迹连续; 不从开球光滑推有限最大值.
\end{solution}
\begin{exercise}\label{ps:5III}
\textbf{原题 III.} $u$ 在 $\R^3$ 调和, $|u(x)|\le\log(1+|x|)$, 证明 $u\equiv0$.
\end{exercise}
\begin{solution}
$u(0)=0$. 在 $B_R$ 中 $v_R=u+\log(1+R)\ge0$ 且调和. 对固定 $x$, $r=|x|<R$, 球 Harnack 给
\[
\frac{R(R-r)}{(R+r)^2}v_R(0)\le v_R(x)\le
\frac{R(R+r)}{(R-r)^2}v_R(0).
\]
两个系数均为 $1+O(r/R)$, 故 $|u(x)|\le C_x\log(1+R)/R$ 当 $R$ 足够大. 令 $R\to\infty$ 得0. 原题 IV、V 的完整商业书题面未公开, 本册不依据提示猜题.
\end{solution}
\originalsection{作业6: 2011年10月20日交}
\begin{exercise}\label{ps:6I}
\textbf{原题 I(a,b).} 对 $u_{tt}=\Delta u$, 置 $J^0=(u_t^2+|\nabla u|^2)/2$, $J^i=-u_tu_i$. 证明 $\partial_\mu J^\mu=0$, 且对欧氏点积、$V=(1,\omega)$, $|\omega|\le1$, 有 $V\cdot J\ge0$.
\end{exercise}
\begin{solution}
(a) 求导后 $\partial_tJ^0+\partial_iJ^i=u_t(u_{tt}-\Delta u)=0$, 混合导数项抵消.
(b) $V\cdot J=(u_t^2+|\nabla u|^2)/2-u_t\omega\cdot\nabla u\ge(|u_t|-|\nabla u|)^2/2\ge0$, 用 $|\omega\cdot\nabla u|\le|\nabla u|$.
\end{solution}
\begin{exercise}\label{ps:6II}
\textbf{原题 II(a--d).} 在后向截锥 $0\le\tau\le t\le R$, $|y-p|\le R-\tau$, 求侧面外单位法向, 证明局部能量 $\int_{B_{R-t}(p)}J^0(t)\le\int_{B_R(p)}J^0(0)$; 推出紧支撑数据的有限传播, 再补论证全空间能量等号.
\end{exercise}
\begin{solution}
(a) 侧面是 $h=\tau+|y-p|-R=0$, 内部 $h\le0$, $\nabla_{\tau,y}h=(1,\omega)$, $\omega=(y-p)/|y-p|$. 因长度为 $\sqrt2$, 外法向 $N=(1,\omega)/\sqrt2$.
(b) 时空散度定理给 $0=E_B(t)-E_B(0)+\int_{\rm side}N\cdot J$. 后项非负, 即局部不等式.
(c) 对 $|x-p|>R_0+t$, 选足够小 $\delta>0$ 使 $B_{t+\delta}(x)$ 与初值支撑不交. 局部估计给在 $B_\delta(x)$ 的 $u_t,\nabla u$ 均为0. 亦可将锥缩至任意 $(s,x)$, $0\le s\le t$, 得 $u_t(s,x)=0$. 因 $u(0,x)=0$, 沿竖直线积分得 $u(t,x)=0$. 这一步排除“梯度零但值为非零常数”的漏洞.
(d) 固定有限时间 $T$, 支撑在 $B_{R_0+T}(p)$. 选固定更大的球, 侧面通量为0, 散度定理给 $\int J^0(t)=\int J^0(0)$. 故 $\|\nabla_{t,x}u(t)\|_2=\bigl(\int(|g|^2+|\nabla f|^2)\bigr)^{1/2}$. 仅令局部不等式的 $R\to\infty$ 只得“$\le$”, 必须另用零通量或时间反演补等号.
\end{solution}
\begin{exercise}\label{ps:6III}
\textbf{原题 III.} 一维紧支撑光滑初值于 $[-R,R]$, 证明当 $t\ge R$ 时 $P^2=\int u_x^2$, $K^2=\int u_t^2$ 均等于 $(P^2+K^2)/2$.
\end{exercise}
\begin{solution}
由 d'Alembert 公式,
$u_x=\tfrac12[f'(x+t)+f'(x-t)+g(x+t)-g(x-t)]$,
$u_t=\tfrac12[f'(x+t)-f'(x-t)+g(x+t)+g(x-t)]$.
置 $A=f'(x+t)+g(x+t)$, $B=f'(x-t)-g(x-t)$, 则 $u_x=(A+B)/2$, $u_t=(A-B)/2$, 所以 $P^2-K^2=\int AB$. $A$ 支撑于 $[-R-t,R-t]$, $B$ 于 $[-R+t,R+t]$, $t\ge R$ 时除单点外不交, 积分为0. 总能量由前题守恒.
\end{solution}
